\src{173}
Rappelons enfin qu'on a aussi une structure multiplicative entre foncteurs
$\symup{H^{p}(X,A)}$ (cup-produit). On constate alors que les suites spectrales de la
proposition 1 sont compatibles avec les structures multiplicatives, de façon
plus précise, on a un accouplement de la suite spectrale $\symup{E(A,B)}$ et de la suite
spectrale $\symup{E(B,C)}$ vers la suite spectrale $\symup{E(A,C)}$, qui, pour l'aboutissement
se réduit à l'accouplement entre les $\symup{\operatorname{Ext}}$ globaux, et pour le
terme initial à l'accouplement déduit, dans le deuxième membre de (1.4), du
cup-produit et des accouplements des $\symup{\operatorname{Ext}}$ locaux. Il en résulte
en particulier que les \fquote{edge-homomorphisms}
\begin{equation}\tag{2.6}
  \symup{\operatorname{Ext}_{\Osh}^{n}(X;A,B)
  \longrightarrow H^{0}\bigl(X;\Extsh_{\Osh}^{n}(A,B)\bigr)}
\end{equation}
\begin{equation}\tag{2.7}
  \symup{H^{n}\bigl(X,\Homsh_{\Osh}(A,B)\bigr)
  \longrightarrow\operatorname{Ext}_{\Osh}^{n}(X;A,B)}
\end{equation}
sont compatibles avec les structures multiplicatives. Si on se borne donc aux
faisceaux localement isomorphes à un $\symup{\Osh^{m}}$, cela explicite complètement la
composition entre $\symup{\operatorname{Ext}}$ globaux à l'aide du cup-produit, compte
tenu de l'isomorphisme (1.5).

\fsection{\normalfont3}{\normalfont\uline{Résultats de cohomologie locale}.}

Soit $\symup{A}$ un anneau commutatif avec unité muni d'un idéal $\symup{J}$. Nous allons
définir, pour un $\symup{A}$-module $\symup{M}$ variable, des homomorphismes fonctoriels
\begin{equation}\tag{3.1}
\left\{
\begin{aligned}\symup{\operatorname{Ext}_{A}^{p}(A/J,M)}&\symup{\longrightarrow
  \operatorname{Hom}_{A}\!\left(\bigwedge^{p}J/J^{2},M\otimes A/J\right)}\\\symup{\operatorname{Tor}_{p}^{A}(A/J,M)}&\symup{\longleftarrow
  \left(\bigwedge^{p}J/J^{2}\right)
  \otimes\operatorname{Hom}_{A}(A/J,M)}\end{aligned}
\right.
\end{equation}
où les produits tensoriels et extérieurs sont pris sur l'anneau $\symup{A}$ ; noter
d'ailleurs que $\symup{J/J^{2}}$ est en fait un $\symup{A/J}$-module, et que ses puissances
extérieures en tant que $\symup{A}$-module ou en tant que $\symup{A/J}$-module sont les mêmes.
La définition des homomorphismes (3.1) revient à la définition, pour tout
système $\symup{\underline{x}=(x_1,\ldots,x_p)}$ de points de $\symup{J}$ d'homomorphismes
$\symup{\varphi_{\underline{x}}}$ :
\begin{equation}\tag{3.2}
\left\{
\begin{aligned}\symup{\varphi_{\underline{x}}:
  \operatorname{Ext}_{A}^{p}(A/J,M)}&\symup{\longrightarrow M\otimes A/J}\\\symup{\varphi_{\underline{x}}:
  \operatorname{Hom}_{A}(A/J,M)}&\symup{\longrightarrow
  \operatorname{Tor}_{p}^{A}(A/J,M)}\end{aligned}
\right.
\end{equation}
de telle façon que les conditions suivantes soient satisfaites :

